\subsection{A point described by direction and distance}

Choose a point $O$, called the pole, and a directed ray from $O$, called the
polar axis. A point $P$ is described by
\[
\boxed{(r,\theta)},
\]
where $r$ is the distance from $O$ to $P$ and $\theta$ is the directed angle from
the polar axis to the ray $OP$.

\begin{center}
\begin{tikzpicture}[scale=1.0,>=stealth]
    \draw[axis,->] (-0.8,0)--(6.2,0) node[right] {polar axis};
    \coordinate (O) at (0,0);
    \coordinate (P) at ({4.2*cos(38)},{4.2*sin(38)});
    \draw[blue!70!black,line width=1.2pt,->] (O)--(P)
        node[midway,above left] {$r$};
    \draw (1.1,0) arc[start angle=0,end angle=38,radius=1.1];
    \node at (1.35,0.48) {$\theta$};
    \fill[geometricSubsetColor] (O) circle (2.2pt) node[below left] {$O$};
    \fill[geometricSubsetColor] (P) circle (2.2pt) node[above right] {$P$};
\end{tikzpicture}
\end{center}

For this chapter we use $r\geq0$ and $0\leq\theta<2\pi$. The pole is represented
by $r=0$; its angle is geometrically irrelevant.

\begin{interpretationbox}[Why this is adapted to a focus]
If the pole is placed at a focus, then the first coordinate $r$ is already the
focal distance. No square root is needed to recover it from $x$ and $y$.
\end{interpretationbox}

The coordinate curves also change character. In Cartesian coordinates,
$x=\text{constant}$ and $y=\text{constant}$ are lines. In polar coordinates,
\[
\theta=\text{constant}
\]
is a ray, while
\[
r=\text{constant}
\]
is a circle centred at the pole.

\begin{center}
\begin{tikzpicture}[scale=0.82,>=stealth]
    \draw[axis,->] (-4.2,0)--(4.4,0);
    \draw[axis,->] (0,-4.0)--(0,4.2);
    \foreach \rad in {1,2,3} {\draw[gray!45] (0,0) circle (\rad);}
    \foreach \ang in {0,30,...,330} {
        \draw[gray!45] (0,0)--({3.4*cos(\ang)},{3.4*sin(\ang)});
    }
    \draw[blue!70!black,line width=1.2pt] (0,0)--({3.5*cos(35)},{3.5*sin(35)});
    \draw[orange!85!black,line width=1.2pt] (0,0) circle (2.2);
    \node[blue!70!black] at (3.0,2.4) {$\theta=35^\circ$};
    \node[orange!85!black] at (-2.8,-1.6) {$r=2.2$};
\end{tikzpicture}
\end{center}
